%%%%%%%%%%  scan idx 447  |  folio 436  %%%%%%%%%%

\medskip
\noindent
9.6. Estamos ya en condiciones de definir el acoplamiento canónico
anunciado (9.1.2), sin restricción alguna sobre $\ell$, de manera que,
para $\ell\neq p$, esta definición coincida con la dada en 9.2. Para
ello aplicamos las consideraciones de 9.5 a la situación siguiente:

\begin{equation}\tag{9.6.1}
\left\{
\begin{aligned}
P&=T_\ell(A^o)^t=\check{\underline M}{}'_\ell(1),
& Q&=\underline M_\ell,
& R&=T_\ell(A^o)^f/T_\ell(A^o)^t,\\[1.25ex]
F&=T_\ell(A^o)^f,
& E&=D(F')=D\bigl(T_\ell({A'}^o)^f\bigr).
\end{aligned}
\right.
\end{equation}
En la última fórmula, $D$ denota la dualidad de Cartier en el sentido
de los pro-$\ell$-grupos de Barsotti--Tate; es decir, $D(Y)$ es el
sistema \underline{proyectivo} de los $D(Y(n))$, con los morfismos de
transición usuales. La estructura de $F$ como extensión de $R$ por
$P$ es clara. Análogamente,
\begin{equation*}
F'=T_\ell({A'}^o)^f
\end{equation*}
posee la estructura de una extensión de $R'$ por
$P'=T_\ell({A'}^o)^t=\check{\underline M}_\ell(1)$. Aplicando la
dualidad de Cartier, se concluye que $E=D(F')$ está dotado de la
estructura de una extensión de $D(P')=\underline M_\ell$ por $D(R')$.
Ahora bien, $D(R')$ es canónicamente isomorfo a $R$ por 5.5
(demostrado en 7.4.6); esto precisa la estructura de $E$ como extensión
de $Q$ por $R$.

%%%%%%%%%%  scan idx 448  |  folio 437  %%%%%%%%%%

Como extensión mixta $X_\eta$ de $E_\eta$ por $F_\eta$, tomamos
\begin{equation}\tag{9.6.2}
X_\eta=T_\ell(A_K)=X^0_\eta
 \supset X^1_\eta=T_\ell(A_K)^f
 \supset X^2_\eta=T_\ell(A_K)^t
 \supset0.
\end{equation}
El isomorfismo de extensiones de $R_\eta$ por $P_\eta$
\begin{equation*}
X^1_\eta\simeq F_\eta
\end{equation*}
es claro. El isomorfismo de extensiones de $Q_\eta$ por $R_\eta$
\begin{equation*}
X_\eta/X^2_\eta\simeq\sisi{F_\eta}{E_\eta}=D(F'_\eta)
\end{equation*}
se deduce inmediatamente del teorema de ortogonalidad 5.2.

Disponemos, pues, de los datos necesarios para definir un acoplamiento
\sisi{(9.5.6)}{(9.5.5)}, que aquí, puesto que
$Q^*=\underline M'_\ell$, adopta la forma (9.1.2). Su compatibilidad
con la definición de 9.2 se deduce de 9.5.8. Las consideraciones de
5.8 muestran además, bajo reserva de la validez del teorema de Tate
[33] (y por tanto al menos cuando $\operatorname{char}K=0$), que el
acoplamiento considerado puede describirse en términos del único
pro-$\ell$-grupo de Barsotti--Tate $X_\eta=T_\ell(A_K)$ sobre
$\sisi{S}{\eta}$: su filtración de tres escalones y las extensiones
$F$ y $E$ que prolongan $X^1_\eta$ y $X_\eta/X^2_\eta$ pueden
definirse canónicamente en términos de $X_\eta$ solamente.

\medskip
\noindent
\underline{Observación 9.7.} Ya no suponemos necesariamente que $A_K$
tenga reducción semiestable sobre el rasgo henseliano $S$. Sea $K'$
una extensión de Galois finita de $K$ tal que $A_{K'}$ tenga reducción
semiestable sobre la normalización $S'$ de $S$ en $K'$ (3.6). Con las
notaciones de (8.1.8), el descenso de Galois a partir del acoplamiento
de monodromía de $A_{K'}$ define entonces un acoplamiento

%%%%%%%%%%  scan idx 449  |  folio 438  %%%%%%%%%%

\begin{equation}\tag{*}
v_\ell^{K'}:\underline M_K\otimes\underline M'_K
 \longrightarrow\mathbb Z_{\ell K}.
\end{equation}
A priori, este acoplamiento depende de la elección de $K'$, pero se
deduce de (10.4.3) más abajo que el acoplamiento
\begin{equation}\tag{9.7.1}
u_\ell:\underline M_K\otimes\underline M'_K
 \longrightarrow\mathbb Q_{\ell K}
\end{equation}
dado por
\begin{equation*}
u_\ell=\frac{1}{e(S'/S)}\,v_\ell^{K'}
\end{equation*}
es independiente de ella, donde
\begin{equation*}
e(S'/S)=\operatorname{length}(V'/\mathfrak mV').
\end{equation*}
Podemos volver a llamar a (9.7.1) el \underline{acoplamiento de
monodromía} del esquema abeliano $A_K$ sobre $K$ y del número primo
$\ell$ considerado: cuando el propio $A_K$ tiene reducción
semiestable, coincide manifiestamente con la restricción a $\eta$ del
acoplamiento de monodromía (9.1.2). Se deduce de 10.4 más abajo que este
acoplamiento toma siempre, de hecho, valores en $\mathbb Q_K$ y es
independiente de $\ell$.

Cabe preguntarse si toma valores en $\mathbb Z_K$, lo que equivale a
decir que (9.7.1) toma valores en $\mathbb Z_{\ell K}$ para todo
número primo $\ell$. Así ocurre al menos cuando $A_K$ es una curva
elíptica. El único caso que requiere verificación es entonces aquel en
que $\underline M_K\simeq\underline M'_K$ tiene rango $1$, de modo que
$\underline M_K\otimes\underline M'_K$ es canónicamente isomorfo a
$\mathbb Z_K$ y el acoplamiento de monodromía (9.7.1) puede
interpretarse simplemente como un número racional ordinario. Se
deduce de los resultados de Néron [19] que este número es $v(j)$,
donde $j$ es el invariante absoluto de $A_K$ y $v$ es la valoración
normalizada de $K$.
